\section{Letters from Evola to Eliade (II)}

\begin{quotex}
\textbf{Julius Evola} was reestablishing his relationship with \textbf{Mircea Eliade} after the war and his own health and legal problems. It is clear that Evola had been keeping up with Eliade's published works. Knowing Eliade's true interests and background, Evola questions Eliade's failure to quote Rene Guenon, and, by implication, his own work. Beside the obvious desire to avoid professional suicide, Eliade intended his works to be scientific, in the objective sense of being free of political or religious partisanship. It is curious that Evola had never read Georges Dumezil before and it seems to have had little influence on his later work. That is unfortunate.
\end{quotex}


15 Dec 1951

Enough time has passed since I received your last letter and our relation was reestablished after the war. In this period I often heard talk about you and your activity. Naturally I read, with interest, the works that you courteously sent me and I did not miss notifying my friends who could be interested in them and probably they have also written to you.

Recently a newly revised and expanded edition of my \emph{Revolt against the Modern World} appeared and I think it was mentioned in it your \emph{Treatise on the History of Religions}\footnote{Not true.}. But about this, and I say it a little tongue in cheek, one should employ some \emph{Vergeltungen}\footnote{Revenge.}. The fact is striking that your works are so overly concerned to not mention any author who does not strictly belong to the official university literature; in your works, e.g., that lovable good man Pettazzoni\footnote{Italian professor of religion.} is abundantly cited, while not a single word is found about Guenon, and not even other authors whose ideas are much closer to those that permit you to certainly orient yourself in the material that you write about. It stands to reason that this is something that concerns only you, but it would be the chance to ask yourself if, all things considered, if imposing these ``academic'' limitations is a game that is worth the candle. I hope that you will not resent these friendly observations.

I've been told that some difficulties arose for Enaudi\footnote{The Italian publisher.} about your translations due to a ridiculous communist veto. Is it true? If so, will these translations be sold again? (I believe they are the \emph{Treatise} and \emph{Yoga}.) It is possible that in this case I can be useful, even if you are already arranging with sufficiently effective relations, like those around Giuseppe Tucci and over Pettazzoni's objection.

As far as it concerns me, in your letter of last year you were so kind to give me the address of a man of your acquaintance, who, you said, could be useful for marketing the French translations of my books. To tell the truth, I wrote him once, but without receiving a response. Now I would like to make other attempts for the same goal. While some relations were already able to be established with other countries, for France still nothing; Mr. Gallimard had taken himself the initiative to ask Laterza for the rights for two of my books after the war, but I have not heard anything more from him since. I permit myself to ask you, dear sir, if you could and would give me some advice in this regard. I would value especially the translation of Revolt which was already accepted by De Noel), and I wondered if the same publisher of your Treatise could be interested in it. Possibly you could introduce me to Mr. Dumezil, who is one of your friends and must exercise an important role with this publisher. Moreover, I now intend to learn more about the work of this author, whom I have read little, because I've been told that he made some interesting contributions to a line of research in which I myself am particularly interested (warrior initiation).

I've also been told that one of your new books, on Shamanism, was just published; if the publisher is still distributing some samples via post, perhaps you will have the courtesy to provide them with my address. Where is your work on sacred orgies found? I have been gathering some material for an essay on sex magic.

I have returned to Rome for good\footnote{After his convalescence and brief imprisonment.} at my old address. Should you ever come to Italy, I hope that we will be able to meet each other here.


\flrightit{Posted on 2012-12-17 by Cologero}

\begin{center}* * *\end{center}

\begin{footnotesize}
\begin{sffamily}
\texttt{Avery on 2012-12-18 at 18:25 said:}

Guenon has some amusing things to say about the ``academic'' in \emph{Traditional Forms \& Cosmic Cycles}. Unfortunately he does not explicitly explain the main problem with modern humanities studies: writers obsess over making their subject a product of their circumstances, and refuse to admit the influence of anything that could stand above history. In order to fit into that model, Eliade thus made a devilish parody of the Traditionalist thesis, castrating the work of Guenon and Evola by turning it into a psychiatric complex.

\hfill

\texttt{Mihai on 2012-12-19 at 04:17 said:}

Actually, Eliade stood his position quite well inside the academic community, to the extent that he could be considered a trojan horse. Though he never directly admitted a supra-human origin of Tradition (obviously), he did quite a nice job in proving that religion is irreducible to factors such as sociology and psychology- see the opening part of his Traité d'histoire des religions.

His work suffered an obvious decline in quality after he became involved with the University of Chicago, but his works published until then offer a great many insights and introductory keys of understanding to various traditional cosmologies, symbols and rites.

The depth of his work is, of course, limited because of the academical environment, but it is no ``devilish parody''

\hfill

\texttt{Mihai on 2012-12-19 at 04:18 said:}

By the way, Guenon's review from TFACC is by no means negative (though, true, it isn't very enthusiastic either).

\hfill

\texttt{Ea Lord of the Deeps on 2013-09-29 at 19:35 said:}

Evola actually read and used Dumezil original thesis source is in this book:
\textit{http://edicionesnuevarepublica.wordpress.com/2010/01/03/\%C2\%ABel-misterio-hiperboreo-escritos-sobre-los-indoeuropeos-1934-1970\%C2\%BB/}

\hfill

\end{sffamily}
\end{footnotesize}

